\documentclass[aspectratio=169,10pt]{beamer}

\usepackage{booktabs}
\usepackage{array}
\usepackage{amsmath}
\usepackage{siunitx}

\definecolor{projectBlue}{HTML}{275DAD}
\definecolor{projectTeal}{HTML}{0B7A75}
\definecolor{projectOrange}{HTML}{D47A22}
\definecolor{projectGray}{HTML}{52606D}
\definecolor{projectLight}{HTML}{F4F7F9}
\definecolor{projectInk}{HTML}{16202A}

\setbeamertemplate{navigation symbols}{}
\setbeamertemplate{footline}{
  \leavevmode%
  \hbox{%
    \begin{beamercolorbox}[wd=.9\paperwidth,ht=2.5ex,dp=1.2ex,leftskip=1em]{author in head/foot}
    \end{beamercolorbox}%
    \begin{beamercolorbox}[wd=.1\paperwidth,ht=2.5ex,dp=1.2ex,center]{author in head/foot}
      \insertframenumber
    \end{beamercolorbox}%
  }%
}

\setbeamercolor{normal text}{fg=projectInk,bg=white}
\setbeamercolor{frametitle}{fg=projectInk,bg=white}
\setbeamercolor{title}{fg=projectInk}
\setbeamercolor{subtitle}{fg=projectBlue}
\setbeamercolor{structure}{fg=projectBlue}
\setbeamercolor{block title}{fg=white,bg=projectBlue}
\setbeamercolor{block body}{fg=projectInk,bg=projectLight}

\setbeamerfont{title}{size=\LARGE,series=\bfseries}
\setbeamerfont{subtitle}{size=\large,series=\bfseries}
\setbeamerfont{frametitle}{size=\Large,series=\bfseries}
\setbeamerfont{normal text}{size=\normalsize}
\setbeamerfont{itemize/enumerate body}{size=\large}
\setbeamerfont{itemize/enumerate subbody}{size=\normalsize}

\newcommand{\smallgray}[1]{\textcolor{projectGray}{#1}}
\newcommand{\best}[1]{\textbf{\textcolor{projectTeal}{#1}}}

\title{Hard vs. Soft Constraint PINN Project}
\author{Filmon Misgane}
\date{}

\begin{document}

%------------------------------------------------
\begin{frame}[plain]
  \vspace{0.8cm}
  {\usebeamerfont{title}\usebeamercolor[fg]{title}\inserttitle\par}

  \vspace{0.8cm}
  \begin{itemize}
    \item What I have worked on
    \item Current progress
    \item Plan for finishing the project from the blueprint
  \end{itemize}

  \vfill
  \begin{flushright}
    \smallgray{Filmon Misgane}
  \end{flushright}
\end{frame}

%------------------------------------------------
\begin{frame}{Project question}
  \smallgray{The comparison is about accuracy, physical validity, uncertainty, and compute cost.}

  \vspace{0.5cm}
  \begin{block}{Main question}
    Does guaranteeing thermodynamic validity by construction outperform
    penalizing violations, and what does it cost in accuracy, calibration,
    and computational effort?
  \end{block}

  \vspace{0.3cm}
  \begin{table}
    \centering
    \begin{tabular}{p{0.28\linewidth}p{0.58\linewidth}}
      \toprule
      \textbf{Evaluation focus} & \textbf{What we check}                                 \\
      \midrule
      Accuracy                  & Prediction error on measured methane adsorption        \\
      Physical validity         & Bounds, monotonicity, Sips shape, and Arrhenius checks \\
      Extrapolation             & Performance on withheld coal rank and high pressure    \\
      Uncertainty               & Whether prediction intervals match observed errors     \\
      \bottomrule
    \end{tabular}
  \end{table}
\end{frame}

%------------------------------------------------
\begin{frame}{Data pipeline}
  \smallgray{The first milestone was a leakage-safe split with fixed evaluation tiers.}

  \vspace{0.25cm}
  \begin{columns}[T,onlytextwidth]
    \begin{column}{0.66\textwidth}
      \begin{table}
        \centering
        \scriptsize
        \begin{tabular}{lrr}
          \toprule
          \textbf{Split}                 & \textbf{Samples} & \textbf{Experiments} \\
          \midrule
          Full dataset                   & 993              & 114                  \\
          Training pool excluding T2     & 752              & 86                   \\
          Training                       & 472              & 54                   \\
          Validation                     & 132              & 14                   \\
          T1 interpolation               & 148              & 18                   \\
          T2 coal-rank extrapolation     & 241              & 28                   \\
          T3 pressure extrapolation      & 3                & 2                    \\
          \bottomrule
        \end{tabular}
      \end{table}
    \end{column}
    \begin{column}{0.30\textwidth}
      \textbf{Key pipeline choices}
      \begin{itemize}
        \item Experiment-level grouping
        \item Train-only scaling
        \item Raw positive pressure preserved
        \item T2 holds out Subbituminous / Lignite
      \end{itemize}
    \end{column}
  \end{columns}

  \vspace{0.12cm}
  \smallgray{\textbf{Note:} T3 is counted inside T1, so it is not added as an independent split.}
\end{frame}

%------------------------------------------------
\begin{frame}{Blueprint progress}
  \smallgray{The implementation is mostly complete. The remaining work is stronger runs and final writing.}

  \vspace{0.35cm}
  \begin{table}
    \centering
    \begin{tabular}{p{0.20\linewidth}p{0.30\linewidth}p{0.40\linewidth}}
      \toprule
      \textbf{Blueprint part} & \textbf{Status} & \textbf{Output}                              \\
      \midrule
      Steps 1--2              & Completed       & Data pipeline and unified evaluation harness \\
      Steps 3--8              & Implemented     & Soft baseline and hard model variants        \\
      Steps 9--11             & Implemented     & Ablation, UQ calibration, interpretability   \\
      Current work            & In progress     & Results analysis and report refinement       \\
      Next work               & Planned         & Final checking, writing, and packaging       \\
      \bottomrule
    \end{tabular}
  \end{table}
\end{frame}

%------------------------------------------------
\begin{frame}[fragile]{Completed training protocol}
  \smallgray{The full V1 tuning run keeps the intended model size and uses the full cross-validation protocol.}

  \vspace{0.25cm}
  \begin{columns}[T,totalwidth=\linewidth]
    \begin{column}{0.55\linewidth}
      \textbf{\textcolor{projectBlue}{Full protocol now completed}}
      \vspace{0.1cm}
      {\fontsize{6}{7}\selectfont
        \begin{verbatim}
hidden_dims = (512, 1024, 512, 256, 128)
max_epochs = 1000
patience = 100
lambda_p_values = (0.0, 0.01, 0.05, 0.1, 0.2, 0.5, 1.0, 2.0)
seeds = (0, 1, 2)
max_folds = 5
      \end{verbatim}}
    \end{column}
    \begin{column}{0.39\linewidth}
      \textbf{\textcolor{projectBlue}{What this means}}
      \vspace{0.15cm}
      {\small
        \begin{itemize}
          \item 15 validation runs per \(\lambda_p\)
          \item 120 V1 tuning runs in total
          \item Final V1 checkpoint uses \(\lambda_p=0.05\)
        \end{itemize}}
    \end{column}
  \end{columns}

  \vspace{0.25cm}
  \smallgray{The remaining work is interpretation, report polishing, and reproducibility checking.}
\end{frame}

%------------------------------------------------
\begin{frame}{Model variants}
  \smallgray{The models compare soft penalties with hard architectural constraints.}

  \vspace{0.25cm}
  \begin{table}
    \centering
    \scriptsize
    \begin{tabular}{lll}
      \toprule
      \textbf{Variant} & \textbf{Constraint design} & \textbf{Purpose}                                    \\
      \midrule
      V1               & All soft penalties         & Reference model                                     \\
      V3               & Hard bounds                & Force adsorption below learned capacity             \\
      V4               & Hard pressure monotonicity & Force adsorption to increase with pressure          \\
      V5               & Hard Sips gauge            & Build the pressure--adsorption curve into the model \\
      V6               & Hard Van't Hoff / Arrhenius & Constrain temperature dependence                   \\
      V7               & Hard bounds + monotonicity  & Test interaction between two constraints           \\
      V2/V8            & All-hard combined variants  & Test maximum enforcement                           \\
      \bottomrule
    \end{tabular}
  \end{table}

  \vspace{0.2cm}
  \smallgray{V4 pressure isolation follows the blueprint: \texttt{pressure\_MPa}, \texttt{Pr}, and \texttt{PT} are excluded from the free branch.}
\end{frame}

%------------------------------------------------
\begin{frame}{Soft-baseline selection}
  \smallgray{The soft baseline chooses \(\lambda_p\) using validation accuracy and physical validity.}

  \vspace{0.2cm}
  \begin{align*}
    \mathcal{L}_{total}
     & = \mathcal{L}_{data} + \lambda_p \mathcal{L}_{physics}                                      \\
    s_{val}
     & = \mathrm{RMSE}_{val}\left(1 + v_{val}\right)                                               \\
    v_{val}
     & = \max\{\text{\texttt{neg\_rate}},\text{\texttt{sat\_rate}},\text{\texttt{nonmono\_rate}}\}
  \end{align*}

  \vspace{-0.15cm}
  \smallgray{\(s_{val}\) is the selection score. \(v_{val}\) is the largest validation rule violation.}

  \vspace{0.08cm}
  \begin{table}
    \centering
    \scriptsize
    \begin{tabular}{rrrr}
      \toprule
      \(\lambda_p\) & \textbf{Val. score} & \textbf{Val. RMSE} & \textbf{Violation rate} \\
      \midrule
      0.00          & 10.349              & 5.174              & 1.000                   \\
      0.05          & \best{4.916}        & \best{4.897}       & 0.004                   \\
      0.50          & 4.969               & 4.966              & \best{0.001}            \\
      1.00          & 4.941               & 4.922              & 0.004                   \\
      \bottomrule
    \end{tabular}
  \end{table}

  \vspace{0.15cm}
  \smallgray{The full run selects \(\lambda_p=0.05\). It keeps the validation RMSE low while reducing the validation violation rate sharply.}
\end{frame}

%------------------------------------------------
\begin{frame}{Full protocol RMSE results}
  \smallgray{Lower values mean smaller adsorption prediction error.}

  \vspace{0.25cm}
  \begin{table}
    \centering
    \begin{tabular}{lrrr}
      \toprule
      \textbf{Variant}              & \textbf{T1 RMSE} & \textbf{T2 RMSE} & \textbf{T3 RMSE} \\
      \midrule
      V1 soft baseline              & 5.300            & 5.746            & 4.943            \\
      V3 hard bounds                & 5.624            & 5.903            & 2.217            \\
      V4 hard monotonicity          & \best{5.213}     & 5.709            & 4.643            \\
      V5 hard Sips gauge            & 5.410            & 5.476            & 1.372            \\
      V6 hard Van't Hoff / Arrhenius & 6.082           & 5.543            & 5.478            \\
      V7 hard bounds + monotonicity & 7.116            & 6.145            & \best{1.192}     \\
      V2/V8 all-hard combined       & 6.209            & \best{5.056}     & 1.435            \\
      \bottomrule
    \end{tabular}
  \end{table}

  \vspace{0.2cm}
  \smallgray{Current pattern: V4 is best on T1, V2/V8 is best on T2, and V7 is best on T3. The V1 row uses the selected \(\lambda_p=0.05\) model. T3 has only 3 samples.}
\end{frame}

%------------------------------------------------
\begin{frame}{Physical validity and trade-offs}
  \smallgray{Hard constraints remove selected violations by construction.}

  \vspace{0.25cm}
  \begin{table}
    \centering
    \scriptsize
    \begin{tabular}{p{0.24\linewidth}p{0.28\linewidth}p{0.38\linewidth}}
      \toprule
      \textbf{Check}  & \textbf{Current result}                      & \textbf{Interpretation}                             \\
      \midrule
      V4 monotonicity & Minimum slope 0.111; violation rate 0        & Pressure isolation matches the blueprint            \\
      V5 Sips gauge   & Predictions stayed physically valid          & Pressure--adsorption shape is built in              \\
      V6 Arrhenius    & Identity error \(1.82\times 10^{-12}\)       & Temperature relation holds numerically              \\
      V2/V8 all hard  & Sips checks pass                             & All checked Sips rules pass                         \\
      \bottomrule
    \end{tabular}
  \end{table}

  \vspace{0.25cm}
  \begin{block}{Trade-off to discuss}
    V4 satisfies the monotonicity check and performs best on T1, but no hard
    variant wins every tier. This supports a careful conclusion rather than a
    broad claim that hard constraints always improve accuracy.
  \end{block}
\end{frame}

%------------------------------------------------
\begin{frame}{Uncertainty calibration}
  \smallgray{Lower ECE means the predicted uncertainty matches the observed errors better.}

  \vspace{0.25cm}
  \begin{table}
    \centering
    \scriptsize
    \begin{tabular}{lrrrr}
      \toprule
      \textbf{Variant}              & \textbf{Val. ECE} & \textbf{T1 ECE} & \textbf{T2 ECE} & \textbf{T3 ECE} \\
      \midrule
      V1 soft baseline              & 0.032             & 0.027           & 0.273           & 0.120           \\
      V3 hard bounds                & 0.042             & 0.053           & 0.242           & 0.120           \\
      V4 hard monotonicity          & 0.039             & 0.026           & 0.321           & 0.120           \\
      V5 hard Sips gauge            & 0.041             & 0.032           & 0.276           & 0.120           \\
      V6 hard Van't Hoff / Arrhenius & 0.045            & 0.036           & 0.308           & 0.120           \\
      V7 hard bounds + monotonicity & \best{0.017}      & 0.015           & 0.299           & 0.120           \\
      V2/V8 all-hard combined       & 0.026             & \best{0.015}    & \best{0.202}    & 0.120           \\
      \bottomrule
    \end{tabular}
  \end{table}

  \vspace{0.2cm}
  {\small
    \begin{itemize}
      \item Uncertainty estimates are less reliable on T2. V2/V8 has the lowest T2 ECE in the full run.
      \item T3 ECE is included for completeness, but it has only 3 samples.
    \end{itemize}}
\end{frame}

%------------------------------------------------
\begin{frame}{Conclusion based on the full protocol}
  \begin{block}{Current answer}
    Hard constraints help most clearly with physical validity. Their accuracy
    and uncertainty benefits depend on the model variant and evaluation tier.
  \end{block}

  \vspace{0.25cm}
  {\small
    \begin{itemize}
      \item Physical validity: hard variants remove the targeted violations by design.
      \item Accuracy: V4 is best on T1, V2/V8 is best on T2, and V7 is best on T3.
      \item Calibration: V2/V8 has the lowest T2 ECE, while V7 has the lowest validation ECE.
            %\item Cost: hard variants need more specialized architecture and stronger verification.
      \item Next step: finalize the report discussion and reproducibility package.
    \end{itemize}}
\end{frame}

%------------------------------------------------
\begin{frame}{Plan for the rest of the project}
  \smallgray{The main remaining work is writing, checking, and packaging.}

  \vspace{0.25cm}
  \begin{table}
    \centering
    \scriptsize
    \begin{tabular}{p{0.22\linewidth}p{0.34\linewidth}p{0.34\linewidth}}
      \toprule
      \textbf{Work}         & \textbf{Why it matters}                      & \textbf{Expected progress}                         \\
      \midrule
      Verify outputs      & Confirms tables match saved CSV files       & Mostly complete after the full run          \\
      Finish report draft & Summarizes methods, results, and limits     & Results discussion needs final polishing    \\
      Supervisor review   & Checks scientific interpretation            & Useful before final manuscript writing      \\
      Final package       & Makes the project reproducible              & Clean results, README, figures, and report  \\
      \bottomrule
    \end{tabular}
  \end{table}
\end{frame}

\end{document}
